\section{Security Analysis}
\label{sec:security}

We analyze the public-coin interactive protocol.  The root \(\rt_{ABC}\)
fixes three arbitrary interleaved words before the scalar challenges are
sampled; we do not assume that these words are valid encodings.  After
\((r,s)\), the prover fixes the intermediate-vector root before the query set
\(I\) and the code-check points are sampled.

For \(W\in\mathbb{F}^{k\times n}\), let
\(\Delta(W,\mathcal{C}^{\langle k\rangle})\) be its relative column distance
from the interleaved code.  We use the following code property; the full
parameter discussion is given in Appendix~\ref{app:security-proof}.

\begin{definition}[Structured-fold proximity gap]
\label{def:structured-proximity-gap}
For \(\vect{a}=(1,a,\ldots,a^{k-1})\), define
\(\mathsf{FoldWord}(\vect{a},W)[j]=\fold(\vect{a},W[*,j])\).  The code has a
\((k,\tau,\epsilon_{\mathsf{gap}})\) structured-fold proximity gap if every
\(W\) at distance greater than \(\tau\) folds to a word within distance
\(\tau\) of \(\mathcal{C}\) with probability at most
\(\epsilon_{\mathsf{gap}}\) over \(a\sample\mathbb{F}\).
\end{definition}

\begin{definition}[Backend assumptions]
\label{def:lamp-backend-assumptions}
The CP-SNARK, sampled openings, and CP-Link are sound with errors
\(\epsilon_{\mathsf{cp}}\), \(\epsilon_{\mathsf{open}}\), and
\(\epsilon_{\mathsf{link}}\), respectively.  The four intermediate-vector
encoding checks have error \(\epsilon_{\mathsf{code}}\); for the
Reed--Solomon check of Appendix~\ref{app:rs-barycentric},
\[
  \epsilon_{\mathsf{code}}
  \le \frac{4(n-1)}{|\mathbb{F}|-n}.
\]
For zero knowledge, the CP-SNARK and CP-Link are zero knowledge and the
Pedersen commitments are hiding.
\end{definition}

\begin{theorem}[Completeness, soundness, and zero knowledge]
\label{thm:security}
Let \(\mathcal{C}\) have relative distance \(\delta\), fix
\(0<\tau<\delta/2\), and assume Definition~\ref{def:structured-proximity-gap}
and Definition~\ref{def:lamp-backend-assumptions}.  If \(\rt_{ABC}\) is fixed
before \((r,s)\), then the interactive \(\protocol\) protocol has perfect
completeness and zero knowledge.  Moreover, for any PPT prover, let
\(\mathsf{BadInput}\) denote the event that an input word is farther than
\(\tau\) from the interleaved code or that the uniquely decoded matrices
satisfy \(\mat{A}\mat{B}\ne\mat{C}\).  Then
\[
\begin{aligned}
\Pr[\mathsf{Acc}\land\mathsf{BadInput}]
&\le \epsilon_{\mathsf{cp}}+\epsilon_{\mathsf{open}}
   +\epsilon_{\mathsf{link}}\\
&\quad +\epsilon_{\mathsf{code}}+3\epsilon_{\mathsf{gap}}
   +3(1-\tau)^t\\
&\quad +\frac{k-1}{|\mathbb{F}|}
   +4(1-\delta+\tau)^t .
\end{aligned}
\]
\end{theorem}

\begin{proof}[Proof sketch]
Condition on the soundness of the CP-SNARK, openings, CP-Link, and
intermediate encoding checks.  The roots then bind all words before \(I\), and
the circuit uses their authenticated sampled symbols.

The \(r\)-folds test the input words for \(\mat{A}\) and \(\mat{C}\).  The
coefficient vector \(\vect{x}=\vect{r}\mat{A}\) need not be uniform, so the
\(\vect{x}\)-weighted product equation alone does not give the same guarantee
for \(\mat{B}\); the independent \(s\)-fold supplies that test.  By the
structured-fold gap, any input farther than \(\tau\) from the code is rejected
except for \(3\epsilon_{\mathsf{gap}}+3(1-\tau)^t\).

Otherwise, \(\tau<\delta/2\) gives unique decoded matrices.  If their product
relation is false, structured Freivalds fails with probability at most
\((k-1)/|\mathbb{F}|\).  Outside that event, at least one sampled relation
disagrees on a \((\delta-\tau)\)-fraction of positions, contributing
\(4(1-\delta+\tau)^t\).  Zero knowledge follows from the simulators of the
CP-SNARK and CP-Link and the hiding of Pedersen commitments.  The complete
bad-event decomposition appears in Appendix~\ref{app:security-proof}.
\end{proof}

\begin{remark}[Fiat--Shamir]
The implementation derives \((r,s)\), \(I\), and the code-check points from
domain-separated transcript hashes.  A formal analysis of this multi-round
Fiat--Shamir transform is outside the interactive theorem and left to future
work.
\end{remark}
